% FILE: 01_research_question.tex
% Project: ggen — Governance Generation Engine

\section{Research Question}
\label{sec:ggen-rq}

\subsection{Problem Statement}
\label{sec:ggen-rq-problem}

Process intelligence research produces conformance verdicts, witness lattice positions, and
lifecycle state machines. These outputs are valuable within a laboratory context but are
disconnected from the capital-market and autonomic-runtime consumers that must ultimately act on
them. A board of directors cannot read a SPARQL result set. An autonomic orchestrator cannot
execute an ontology triple. A TypeScript visualizer cannot consume a Petri net proof directly.

The problem ggen addresses is: \emph{how can process intelligence findings be reliably
materialized into the artifact formats required by capital-market consumers and autonomic runtimes,
while preserving an unbroken, cryptographically verifiable chain of provenance from the underlying
ontological evidence to the board-ready deliverable?}

This problem has three dimensions. First, the \emph{format gap}: conformance verdicts must be
rendered into PowerPoint slides, Excel workbooks, and Rust source code without loss of the
evidentiary basis. Second, the \emph{admissibility gap}: capital-market artifacts must satisfy
explicit thresholds---fitness $\geq 0.95$, precision $\geq 0.90$, BLAKE3 receipt chain
intact---before they can be presented to a board. Third, the \emph{repeatability gap}: any claim
in the deck must be independently verifiable by re-executing the originating SPARQL query against
the ontology and comparing the computed fitness to the slide assertion.

\subsection{Old-Regime Assumption Challenged}
\label{sec:ggen-rq-old-regime}

The old-regime assumption is that capital-market artifacts---acquisition decks, due-diligence
workbooks, integration risk assessments---are prepared by human analysts who translate technical
findings into prose and slides. Under this assumption, the chain of custody from measurement to
board presentation is mediated by human judgment, subject to selection bias, and not independently
verifiable.

A secondary old-regime assumption concerns code generation: that autonomic runtime code is
hand-authored by engineers who read specification documents and translate their intent into Rust
structs and methods. Under this assumption, the correspondence between a governance specification
and the runtime behavior it is supposed to enforce is a matter of discipline rather than
mechanical derivation.

ggen challenges both assumptions simultaneously. It posits that capital-market artifacts and
autonomic runtime code can both be \emph{manufactured} from a shared ontological source through
a deterministic, receipt-sealed pipeline. Under this posture, the board deck is not a human
translation of a conformance verdict; it is a Tera template rendered against a SPARQL result set
whose inputs and outputs are both hash-chained. The claim on the slide and the query that produced
it are the same artifact, separated only by the rendering step.

\subsection{Hypothesis}
\label{sec:ggen-rq-hypothesis}

The hypothesis is that a configuration-driven manufacturing pipeline---in which SPARQL queries
extract typed facts from an RDF knowledge base and Tera templates render those facts into output
artifacts, with each output sealed by a BLAKE3 cryptographic receipt---can satisfy three
conditions simultaneously:

\begin{enumerate}
  \item \textbf{Format completeness}: The pipeline can manufacture artifacts in all formats
    required by capital-market consumers (PowerPoint JSON, Excel JSON, Rust source, WIT interface
    definitions, TypeScript bindings) from a single ontological source.
  \item \textbf{Admissibility enforcement}: The pipeline can enforce board admissibility
    thresholds (fitness $\geq 0.95$, precision $\geq 0.90$, log format \texttt{ocel:2.0} or
    \texttt{xes:1849-2016}) as SPARQL \texttt{FILTER} conditions, ensuring that no
    non-conforming claim reaches the output artifact.
  \item \textbf{Provenance completeness}: Every output artifact can be traced back to its
    originating SPARQL query, ontology facts, and conformance verdicts through a BLAKE3
    hash-chained receipt sequence, enabling independent verification by any party with access to
    the ontology.
\end{enumerate}

\subsection{Success Criteria}
\label{sec:ggen-rq-success}

Success requires all three conditions to hold simultaneously, with receipt evidence for each:

\begin{itemize}
  \item The manufacturing pipeline executes without error across all active generation rules, with
    \texttt{validation\_passed: true} recorded in \texttt{audit.json}.
  \item The BLAKE3 receipt chain is unbroken: each receipt records the hash of its predecessor,
    from the chain root (operation \texttt{de2ad19f}, \texttt{previous\_receipt\_hash: null})
    through the final receipt (operation \texttt{09084e1e},
    \texttt{previous\_receipt\_hash: 057521e5\ldots}).
  \item All 42 or more proof gates in the seven audit scripts pass, covering feature law,
    TypeScript projection surface, monomorphization, brand tokens, enum tagging, WASM boundary
    isolation, and component boundary separation.
  \item The embedded test suites in the manufactured Rust artifacts pass: \texttt{test\_witness\_lattice\_monotonicity},
    \texttt{test\_witness\_partial\_order}, \texttt{test\_orchestrator\_initialization}, and
    \texttt{test\_mape\_k\_cycle}.
  \item The visualizer-dashboard-nextjs generation rule is recovered from its disabled state and
    produces a deployable TypeScript artifact, closing the Phase 6 recovery gap documented in
    \texttt{ggen.toml}.
\end{itemize}

The current status is PARTIAL: the receipt chain is intact, the audit.json records
\texttt{validation\_passed: true}, and the Blue River orchestrator artifact is hash-sealed, but
the later receipts record \texttt{output\_hashes: []} indicating that the second and third
manufacturing runs did not produce new output artifacts. The visualizer rule remains disabled.
Full ALIVE status requires resolving both of these gaps.
